\input{1-preamble}

\title{Coercive Threats, Allied Commitment, and Public Support for Defiance}

\author{Hannah Jakob Barrett}

\begin{document}
\maketitle
\begin{center}
    Pre-Analysis Plan
\end{center}

%-------------------------------------------------------------------
\section{Introduction}
%-------------------------------------------------------------------
Coercion is a widely studied phenomenon in international relations. In conflict scenarios, coercive threats are designed to elicit a state's desired outcomes by dangling the threat of harm, which manipulates the risk associated with defiance. \textcite{schelling_arms_1966} originally theorized that coercive threats can be either deterrent or compellent; deterrent threats aim to stop an adversary from acting, while compellent threats aim to force them to change course. Moreover, \textcite{schelling_arms_1966} explicitly theorizes that it is much harder to coerce a state into action, thus changing the status quo, than it is to deter an action and preserve the status quo. 

Subsequent research has corroborated Schelling's claim that coercion suffers from low success rates. Some explanations for coercion failure are structural. Based on the relative power capabilities of the states involved, and the terms of the ultimatum, the credibility of both the initial threat and subsequent assurances may vary in effectiveness. \autocite{fearon_domestic_1994, schelling_arms_1966} Some states face the calculus of resisting coercion in order to set a precedent of resolve and avoid subsequent threats in the future \autocite{sechser_goliaths_2010, sechser_reputations_2018}. This theoretical account also argues that both leaders and citizens are concerned about their nation's reputation for resolve \autocite{yarhi-milo_tying_2018, lupton_reputation_2020, dafoe_reputations_2021} By requiring the target state to publicly accept a loss, coercive threats can activate concerns about national honor, causing even greater public support for defiance and escalation.\autocite{dafoe_coercion_2021, dafoe_provocation_2022} If citizens categorically oppose compliance with coercive threats, democratic leaders risk facing domestic audience costs if they do not stand firm, even when the strategic calculus might favor compliance.\autocite{fearon_domestic_1994, tomz_domestic_2007, davies_audience_2013, levy_backing_2015, yarhi-milo_tying_2018} On the other hand, coercive threats might be written off as cheap talk or bluster.\autocite{fearon_domestic_1994, snyder_cost_2011, yarhi-milo_tying_2018} Insights from political psychology have found that coercion prompts anger, as well as a psychological motivation to resist. More specifically, when coercive threats are perceived to limit the target state's freedom, reactance results in greater support for defiance.\autocite{schaub_deterrence_2004, powers_psychology_2023}  

I argue that the experimental literature on coercion suffers from limitations that systematically bias it toward finding evidence of failure. The first limitation is one of sample selection. The majority of survey experiments on coercion have been conducted on American respondents. \footnote{\citeauthor{dafoe_provocation_2022}'s study on China is an exception.} Studying the effects of coercive threats on public audiences in military superpowers has clear consequences for both theorizing the mechanisms of coercion and finding generalizable evidence of how coercion operates in the real world. With relative power on their side, respondents in strong states are more likely to interpret adversary threats as bluffs, more likely to abhor reputational costs, and more likely to react defiantly. When the adversary's threats are more likely to be dismissed, compliance with coercive threats is the least likely outcome by design. Generalizing findings from US samples to a broader population of states is therefore problematic, given that most states in the international system are not military superpowers. 

%------------------------------- Contribution
This pre-analysis plan addresses this gap by studying public responses to a coercive threat in Denmark. Denmark is a small NATO member with direct exposure to Russian military pressure in both the Baltic and Arctic regions. By situating the adversary as a clearly more powerful state (Russia), the design tests whether the coercive failure mechanisms documented in US samples still hold when the costs of defiance are objectively higher. This study makes an additional novel contribution by introducing a previously unexplored variable to this literature: the role of allied support. For states that lack relative military advantages, the availability and commitment of allied support is essential for responding to coercive threats. Allied support has been shown to deter threats from adversaries, and the presence of military backing may embolden target states to retaliate against adversaries. \autocite{leeds_alliances_2003, fang_concede_2014, benson_inducing_2014} Conversely, uncertain commitments from allies may decrease public support for defiance by raising the costs that the state must bear on its own. Whether and how allied backing shapes the perceived costs of compliance versus defiance is fundamental to understanding how alliances function in crisis situations. Yet direct experimental tests of this mechanism remain rare, and no study has examined this emboldening effect at the level of public opinion. In the current geopolitical context, where the credibility of US security commitments to European allies has come under scrutiny, this question is both timely and theoretically novel.  

%-------------------------------------------------------------------
\section{Theory}
%-------------------------------------------------------------------

\subsection{Reputational Costs of Compliance}
In international conflicts or negotiations, a coercive threat creates a risk calculus for the target state. Whereas compliance with the threat requires capitulation, defiance incurs the risk of military escalation.\autocite{schelling_arms_1966} This risk calculus partially explains why compellence is more likely to fail than deterrence.\autocite{schelling_arms_1966} When compliance with an ultimatum is public, it also involves reputational costs. A state that yields to coercive pressure signals to the adversary and a broader international audience that further threats will also be effective. 

Resistance may therefore be rational even for small states, because compliance sets the precedent that the target state is vulnerable to coercion. \textcite{sechser_goliaths_2010} shows that even strong military powers fail to extract demands from smaller states, because targets resist accepting the downstream costs of visible capitulation. This dynamic operates at the level of public opinion as well. \textcite{dafoe_coercion_2021, dafoe_provocation_2022} demonstrate experimentally that coercive threats activate honor and reputational concerns among ordinary citizens, increasing public support for defiance and escalation as a means of saving face.

The current wisdom in the experimental literature suggests that this reputation mechanism produces a public preference for defiance in response to coercive threats. Citizens who learn that their country has been provoked, issued an ultimatum, and threatened with military force tend to prefer standing firm. The key question is whether this insight travels outside of the United States and China, where relative power advantages may set the perfect scope conditions for defiance. For smaller states on the receiving end of coercive threats, the presence or absence of allied commitments are more consequential for the perceived costs of defiance. 

\subsection{Reactance to Coercive Threats}
A second pathway to defiance operates through emotional mechanisms. Building on reactance theory from psychology, \textcite{powers_psychology_2023} show that coercive threats increase public support for defiance by framing the adversary's demand as an infringement on the target state's freedom. When citizens perceive their country is being told what to do under threat of violence, they instinctively resist and become angry. This emotional mechanism complements the reputational mechanism described above. Reputational concerns predict defiance because compliance is costly; \autocite{dafoe_coercion_2021, dafoe_provocation_2022} reactance predicts defiance because coercion is resented.\autocite{powers_psychology_2023} 

Studies from this literature have also found that weak threats are perceived as less credible than stronger, more explicit threats \autocite{lupton_threat_2024}. \textcite{powers_psychology_2023} also acknowledge that coercive threats may \textbf{suppress} public support for defiance if the public perceives the adversary's threat to be a costly signal of resolve. This has yet to be tested in a literature dominated by survey research on US samples, where in the eyes of the American public, global adversaries are more likely to appear weaker and less capable than the US military.\autocite{pomeroy_hawks_2024} 

If the public perceives coercive threats from a larger power to be a credible signal of resolve, we should expect support for defiance to be suppressed. I argue that when a target state lacks allied backing, the risks associated with defiance may be high enough to suppress the face-saving and reactant mechanisms. In fact, this study is designed to test how these factors interact to shape how the public perceives the risk of defiance in coercive, interstate disputes. The following section elaborates the ways in which alliances may embolden targets to retaliate, and the conditions under which that effect may be mitigated.  

\subsection{Allied Support as a Moderator of Defiance}
Military alliances can affect the costs of defiance in two theoretically distinct ways: through an \textit{emboldening} effect and a \textit{restraining} effect.\autocite{fang_concede_2014} The emboldening effect suggests that allied military backing reduces the perceived risks of defiance. A target state with allies committed to its defense can stand firm against a coercive threat at lower cost, since the risks of military escalation are distributed across the alliance. This risk-sharing function of alliances is central to their deterrent value,\autocite{leeds_alliances_2003, benson_inducing_2014} but may also be activated in compellent threat scenarios. 

This study aims to test whether this effect also operates at the level of public opinion. Alliances have been shown to have a positive effect on public support for military intervention, when an ally faces external threats \autocite{tomz_military_2021, tomz_how_2023} Alliances do this by creating formal commitments, such as mutual defense agreements, and by establishing a moral relationship between states, in which concerns for loyalty drive public willingness to support allies. If emboldening plays out at the public level, citizens who believe their country's allies will come to their defense should update their perceptions about the risks of defiance, and become more willing to stand firm.

The restraining effect operates in the opposite direction.\autocite{fang_concede_2014} Allies may recommend that the target state comply with an adversary's demand in order to avoid escalation, preserve alliance unity, or protect broader strategic interests. Deviating from the preferences of one's allies can create additional \textit{alliance reputation} costs, wherein ignoring allied counsel and retaliating recklessly communicates that one is an unreliable partner.\autocite{fang_concede_2014} Reckless escalation that alienates committed allies may thus be perceived as risky even among citizens who are predisposed toward standing firm. If allies recommend compliance, the public may reduce their support for defiance out of concern for the country's standing within the alliance. 


\subsection{Perceptions of Alliance Reliability}
\textcite{johnson_alliance_2023} argue that alliances are formed \textquote{under a veil of ignorance over what wars will occur in the future.} What is deemed a reliable alliance today might suffer from credibility and legitimacy issues in the future. Moreover, the prospect of war itself may affect the willingness of actors to bear risks on behalf of their allies. In sum, the reliability of alliances fluctuates over time, and is a product of both power asymmetry and changing political circumstances. \autocite{leeds_reliability_2003, gartzke_why_2004} 

Perceptions of alliance reliability are consequential for the risk calculus of target states on the receiving end of coercive threats. \textcite{schelling_arms_1966} reflected on the implications of differing perceptions of alliance reliability in the context of the Cold War. While European allies were preoccupied with the perception that the US might leave them to fight \textquote{their own futile war against the Soviet Union while the United States saved itself,} the United States also feared the possibility that Europe would allow \textquote{the two larger military opponents} to bear the brunt of weapons expenditure and violence.\autocite{schelling_arms_1966}


Studying how alliance reliability shapes public responses to coercive threats is both a theoretical gap and an urgent empirical priority. The experimental literature on coercion and public opinion has concentrated overwhelmingly on great powers, whose capacity to resist threats rests primarily on their own military resources rather than on allied commitments. For small states facing militarily superior adversaries, the reliability of alliances is not a peripheral consideration but a central factor shaping the decision to defy or comply. If a target state anticipates that their allies will prove unreliable, they may condition their strategy and avoid the path of escalation. Empirically, this question could not be more timely. Confidence in the United States to honor its mutual defense commitments and act responsibly in international affairs has declined markedly since the beginning of the second Trump Administration, as public opinion data across NATO member states show.\autocite{wike_us_2025} The result is a historically unusual moment in which both reassurance \textbf{and} the possible withdrawal of American commitments are ecologically valid. Experimental manipulations that would have seemed implausible a decade ago now map directly onto the political reality facing publics in allied states. In short, there is no better moment to study how allied signals shape the micro-foundations of coercion success or failure.

%What measures can we use to tap into perceptions of reliability?
To this end, this study will include pre-treatment measures of international trust to tap into how Danish respondents judge the reliability of other countries in general. \textcite{brewer_international_2004} provides one such measure, which I will adapt for the Danish case. Including this measure as a covariate will help to reach a more precise estimate of treatment effects, given that diffuse international trust should explain some of the variance in how individuals perceive signals from international allies. Moreover, general predispositions of international trust will likely condition the effect of an alliance signal on one's willingness to defy threats. This suggests that international trust may also moderate the treatment effect.   

%secondary outcome of treatment: confidence in NATO allies
While the main interest of this study is investigating how signals from allies affect public willingness to defy international threats, the same signals may also affect individual confidence in one's allies. These effects are unlikely to be unidirectional, given that the experiment will manipulate alliance reliability through the material commitments and rhetoric of NATO allies. Measuring confidence in whether one's allies will come to their country's defense as a secondary treatment outcome will therefore offer exploratory evidence of how allied signals are perceived by different respondents. 

\section{Additional Drivers of Defiance}
\subsection{Individual Predispositions}
%Risk Aversion and Resolve
Following the behavioral revolution in international relations research, the literature provides a strong foundation for studying the micro-foundations of individual willingness to defy threats, accept risks, and escalate conflicts. In fact, given that \textcite{schelling_arms_1966} refers to coercion as a manipulation of risk, studying individual risk tolerance is essential to understanding how individuals may react to a threat that inherently provokes the question of \textquote{what risk is worth taking.}\autocite{schelling_arms_1966}

Including measures of risk tolerance\autocite{kertzer_resolve_2016, vieider_common_2015} as a covariate will help to explain some of the variation in individual preferences for defiance. However, additional predispositions must also be accounted for, given that risk-aversion manifests differently in a diverse population of individuals. 

%How Does this Relate to FP Postures 
While we might assume that individual who are predisposed to avoid risks should exhibit a tendency to avoid escalation and conflict, this would be an oversimplification of the \textbf{types} of decisions that individuals deem risky. In this regard, foreign policy postures provide more informative expectations about how risk-aversion should play out in a coercive threat situation. 

Individuals who score highly on cooperative internationalism are often referred to in the literature as \textquote{doves}, while \textquote{hawks} tend to score higher on militant internationalism.\autocite{wittkopf_foreign_1986, rathbun_taking_2016, brutger_dispositional_2018} For doveish individuals, because escalation is associated with higher risk, backing down from a dispute is deemed the safer option.\autocite{kertzer_resolve_2017, brutger_dispositional_2018} Meanwhile, because hawks associate reputational risks with backing down, risk aversion in a hawk should be associated with increased resolve \autocite{kertzer_resolve_2017, brutger_dispositional_2018}. Hawkish and dovish individuals should also associate very different reputational costs with defying versus complying with a coercive threat. For hawks, giving into a threat is more costly because it involves visible costs to one's national honor and reputation for standing firm \autocite{kertzer_resolve_2017, brutger_dispositional_2018, lupton_reputation_2020, dafoe_reputations_2021} 

Understanding the dispositional differences that cause risk aversion and reputation to lead to different outcomes is fundamental to any study interested in investigating the causal link between threats and defiance. Therefore I include a validated battery of survey items to measure both foreign policy postures and risk tolerance. \autocite{gravelle_structure_2017, kertzer_resolve_2016, vieider_common_2015}

%Attitudes Towards Adversary 

\subsection{Emotions}
Anger has been experimentally associated with greater support for retaliation against external threats,\autocite{shandler_cyber_2021, fisk_emotions_2019} while other emotional reactions such as fear and anxiety have been associated with more cautious, de-escalatory preferences.\autocite{huddy_threat_2005, wagner_anxiety_2019} The experimental literature suggests that because coercive threats are typically perceived as blustery or constraining, they are more likely to activate anger than fear, reinforcing support for defiance through an emotional channel.

The contrasting relationships between anxiety and risk aversion \autocite{huddy_threat_2005, wagner_anxiety_2019} and anger and defiance \autocite{shandler_cyber_2021, fisk_emotions_2019, powers_psychology_2023} thus offer interesting ground for exploratory analyses. 



%-------------------------------------------------------------------
\section{Confirmatory Hypotheses}
%-------------------------------------------------------------------
Insights from the literature on reactance, reputation, and alliance reliability offer distinct predictions about public support for defiance, depending on the nature of the allied signal. This leads to the following pre-registered, confirmatory hypotheses:
\begin{enumerate}
    \item \textbf{H1 (Emboldening)}: When allies express commitment, public support for defiance will increase. 
    \item \textbf{H2 (Restraining)}: When allies express a preference for restraint, public support for defiance will decrease. 
    \item \textbf{H3}: 
\end{enumerate}

For both of these hypotheses, the control group (no allied signal) will serve as the reference. 

\subsection{Exploratory Analyses}
Subgroup effects will be run to compare treatment effects and defiance attitudes across hawks and doves, based on pre-treatment scores on the militant internationalism and cooperative internationalism scales. A similar approach will be taken with the items for risk tolerance,  international trust, and attitudes towards Russia. 

Given the importance of emotions for studying individual support for defiance, an emotions battery will be used to investigate whether threats generate anxiety or anger, and if these emotions are associated with lower or higher support for defiance, respectively. While these directional effects have been demonstrated in other experimental studies, the link between emotions and defiance may be attenuated in a small-state like Denmark, given the risks of escalation. 

Finally, measuring reactance and concerns for reputation across all treatment conditions allows this study to probe how the relationship between reactance, reputation, and support for defiance fluctuates across treatment conditions (which vary whether an alliance commitment is present). 

%-------------------------------------------------------------------
\section{Research Design}
\label{sec:designB}
%-------------------------------------------------------------------

\subsection{Sample and Recruitment}

Participants will be recruited from YouGov's Danish online panel. The target sample is (3000) voting-eligible adults (age 18 and above), allocated equally across treatment conditions (xxx per condition). Respondents who fail attention checks will be excluded from the analytical sample prior to hypothesis testing; target cell sizes refer to the post-exclusion sample.

\subsection{Treatments}
\begin{table}[H]
\centering
\caption{Three Condition Design}
\label{tab:studyb_treatments}
\footnotesize
\begin{tabular}{p{3cm} p{3cm} p{8cm}}
\toprule
\textbf{Condition} & \textbf{Threat} & \textbf{Allied Signal} \\
\midrule
C & Coercive Threat & None \\ 
T1 & Coercive Threat & Denmark's NATO allies reached a consensus and agreed to provide military support to Denmark in this dispute. \\
T2 & Coercive Threat & Denmark's NATO allies urged Denmark to comply with the ultimatum, warning that military escalation would not serve the alliance's collective interests. \\
\bottomrule
\end{tabular}
\end{table}


\begin{table}[H]
  \centering
  \caption{Survey Flow}
  \label{tab:surveyflow}
  \footnotesize
  \begin{tabular}{p{1cm} p{3.5cm} p{8.5cm}}
  \toprule
  \textbf{Block} & \textbf{Content} & \textbf{Measures} \\
  \midrule
  1 & Pre-treatment Block & Age, gender, education, partisanship, political interest, foreign policy postures (isolationism, militant internationalism, cooperative internationalism)
   \\[6pt]
  2 & Attention Check &  \\[6pt]
  3 & Initial Threat Scenario \\ [6pt]
  3 & Manipulation Check  & Did Russia issue a direct demand? \textit{(For All Respondents)} \\[6pt]
  5 & \textit{Ally Treatment} & \textit{(Only for T1 and T2)} \\[6pt]
  6 & \textit{2nd Manipulation Check} & \textit{(Only for T1 and T2: What position did Denmark's allies take?} \\[6pt]
  7 & Outcome Measures & Binary DV (comply/defy), continuous support for escalation (Likert scale), anger battery, reactance battery, reputation battery \\[6pt] 
  8 & Open-ended response & \textit{In a few sentences, please explain why you think Denmark should follow that approach.} \\
  \bottomrule
  \end{tabular}
  \end{table}


%-------------------------------------------------------------------
\section{Questionnaire}
\label{sec:measurement}
%-------------------------------------------------------------------
\subsection{Pre-Treatment Measures}
\begin{enumerate}
    \item \textbf{Demographics:} Age, gender, education level
    \item \textbf{Partisanship:} Prospective vote intention (party list)
    \item \textbf{Political interest:} \textit{How interested are you in politics?} (1 = Not at all interested, 5 = Very interested)
    \item \textbf{Foreign policy postures} Respondents rate their agreement with each item on a 5-point scale from \textit{Strongly disagree} to \textit{Strongly agree}:
    \begin{enumerate}
        \item \textit{Isolationism 1:} Denmark's interests are best protected by avoiding involvement with other nations.
        \item \textit{Isolationism 2:} Denmark needs to simply mind its own business when it comes to international affairs.
        \item \textit{Militant internationalism 1:} Denmark should take all steps including the use of force to prevent aggression by any expansionist power.
        \item \textit{Militant internationalism 2:} Denmark needs a strong military to be effective in international relations.
        \item \textit{Cooperative internationalism 1:} When deciding its foreign policies, Denmark should take into account the views of its major allies.
        \item \textit{Cooperative internationalism 2:} Denmark should be more committed to diplomacy and not so fast to use the military in international crises.
    \end{enumerate}
    \item \textbf{Risk Tolerance}
    \begin{enumerate}
        \item In general, people often think about risks when making financial, career, or other life decisions. Overall, how would you place yourself on the following scale? [1 = very uncomfortable taking risks...3= neither/nor...5 = very comfortable taking risks.]
        \item People behave differently in different situations. How would you rate your willingness to take risks in the following areas? [1 = very uncomfortable taking risks...3= neither/nor...5 = very comfortable taking risks.]
        \begin{enumerate}
            \item While driving
            \item During fitness or leisure
            \item In financial matters
            \item With your own health
            \item In your career
        \end{enumerate}
    \end{enumerate}
    \item \textbf{International Trust:} Generally speaking, how often do you think Denmark can trust other nations? 
    \begin{enumerate}
        \item Always
        \item Most of the time
        \item About half the time
        \item Some of the time
        \item Never
    \end{enumerate}
    \item \textbf{Pre-Treatment Attitudes toward Russia:} 
    \begin{enumerate}
        \item How does Russia's military strength compare to that of Denmark? (On a 5pt scale, 1 = much weaker, 2 = slightly weaker, 3 = equal, 4 = slightly stronger, 5 = much stronger)
        \item \textit{On a scale from 0 to 10, how much of a threat does Russia pose to Denmark's security?}
        \begin{enumerate}
            \item Respondents select one option on a 10-point scale from \textit{Not threatening at all} to \textit{Extremely threatening}
        \end{enumerate}
    \end{enumerate}
\end{enumerate}

\subsection{Vignette}

The crisis scenario is presented as a short news article. 

\bigskip
\noindent\rule{\textwidth}{0.4pt}
\smallskip

\noindent {\small \textsc{International News}}

\medskip

\noindent \textbf{Denmark Detains Russian Vessel}

\medskip

Danish naval forces intercepted a Russian vessel in the Baltic Sea, after it was found shadowing a Danish patrol ship at close range. The Danish navy detained the vessel and crew, which are now under investigation.

\medskip

Russian authorities deny any wrongdoing. In a public statement, the Russian government issued the following demand: if Denmark does not immediately release the vessel and crew, Russia will not hesitate to take military action to secure their release.
\medskip

[One Sentence - Ally Treatment for T1 and T2]

\noindent\rule{\textwidth}{0.4pt}
\bigskip

\subsection{Manipulation Checks}

Following the vignette and prior to the outcome questions, respondents complete the following manipulation checks:

\begin{enumerate}
    \item \textit{First Manipulation Check (All):} In the scenario you just read, did Russia issue a direct demand to Denmark?
    \begin{enumerate}
        \item Yes
        \item No
        \item They requested something but did not demand it
        \item I don't know
    \end{enumerate}
    \item \textit{Second Manipulation Check (T1 and T2):} In the scenario you just read, what position did Denmark's allies take?
    \begin{enumerate}
        \item They expressed support 
        \item They recommended restraint
        \item They stayed silent
        \item I don't know
    \end{enumerate}
\end{enumerate}

Respondents who fail manipulation checks will not be excluded from the primary analysis. A robustness check will replicate the main model test on the subset of respondents who pass the manipulation checks assigned to their treatment group.

\subsection{Primary Outcome: Comply or Defy}

The primary dependent variable confronts respondents with a direct binary choice:

\begin{quote}
\textit{What do you think Denmark should do in response to Russia's ultimatum?}

\quad 1 --- Release the vessel and crew immediately \hfill [Comply]

\quad 2 --- Continue holding the vessel \hfill [Defy]
\end{quote}

\noindent The item is coded 0 (Comply) and 1 (Defy). Following this choice, respondents indicate the strength of their preference:

\begin{quote}
\textit{On a scale from 0 to 10, how strongly do you feel about this choice?}
\end{quote}

\subsection{Defiance Scale}

All respondents are then asked:

\begin{quote}
\textit{If Denmark were to continue holding the vessel, to what extent would you support or oppose each of the following approaches? Please rate each option on a scale from ``Strongly support'' to ``Strongly oppose.''} (5-point scale)
\end{quote}

\begin{enumerate}
    \item Pursue de-escalation through diplomatic channels, without taking any military action
    \item Impose economic sanctions on Russia
    \item Deploy additional Danish naval forces to the area to confront Russia's vessels
    \item Conduct a military strike against Russian naval forces
\end{enumerate}

These four items constitute an ordered defiance scale spanning de-escalation to the use of military force. The scale is qualitatively distinct from the binary dependent variable, in that it measures \textit{how much} defiance the respondent favors and at what cost. The de-escalation item captures respondents who favor standing firm while also minimizing escalation risk. The military strike item captures those willing to pay higher costs of defiance.

\subsection{Open-Ended Response}

Following the escalation scale, respondents are asked:

\begin{quote}
\textit{In a few sentences, please explain why you think Denmark should follow that approach.}
\end{quote}

\subsection{Mechanism Measures}

\paragraph{Emotions battery.}
\textit{Please indicate your feelings toward Russia when thinking about this situation. To what extent do you feel\ldots} (presented in random order; 0 = not at all, 10 = very much)

\begin{enumerate}
    \item Angry \hfill \textit{[anger]}
    \item Furious \hfill \textit{[anger]}
    \item Anxious \hfill \textit{[fear]}
    \item Worried \hfill \textit{[fear]}
    \item Sad \hfill \textit{[distractor]}
    \item Happy \hfill \textit{[distractor]}
\end{enumerate}

An anger index and a fear index are computed as the mean of the respective items.

\paragraph{Reputation battery.}
\textit{To what extent do you agree or disagree with the following statements?} (presented in randomized order; 5-point scale from \textit{Strongly disagree} to \textit{Strongly agree})

\begin{enumerate}
    \item Giving in to Russia puts Denmark's international reputation at stake.
    \item Giving in to Russia would be humiliating for Denmark
    \item Giving in to Russia now would invite further threats to Denmark's security in the long run.
    \item Releasing the vessel would reduce the risk of war with Russia. \hfill \textit{[reverse-coded]}
\end{enumerate}

\paragraph{Reactance battery.}
\textit{To what extent do you agree or disagree with the following statements?} (presented in randomised order; 5-point scale from \textit{Strongly disagree} to \textit{Strongly agree})

\begin{enumerate}
    \item Russia tried to make a decision for Denmark
    \item Russia tried to manipulate Denmark
    \item Russia tried to pressure Denmark 
    \item Russia tried to take away Denmark's freedom to choose how to respond
\end{enumerate}

\subsection{Secondary Outcomes}
\paragraph{Confidence in NATO Allies}
\begin{enumerate}
\item How confident are you that NATO members will come to Denmark's defense?\footnote{Should the scale be ordinal as shown here, or expressed as a percentage?}
\begin{enumerate}
    \item NATO members will definitely not come to Denmark's defense
    \item NATO members will probably not come to Denmark's defense
    \item NATO members may or may not come to Denmark's defense
    \item NATO members will probably come to Denmark's defense
    \item NATO members will definitely come to Denmark's defense
\end{enumerate}
\end{enumerate}
%-------------------------------------------------------------------
  \section{Statistical Models}
  \label{sec:model}
  %-------------------------------------------------------------------

  \subsection{Primary Analyses}

  The primary analysis estimates the average treatment effect of each
  condition relative to the no-commitment control. Informally:
  %
  \[
    \texttt{defiance} \sim \texttt{treatment group}
  \]
  %
  Formally, let $Y_i \in \{0,1\}$ denote the binary comply/defy outcome
  for respondent $i$, where $1 = \text{Defy}$. The estimating equation is:
  %
  \begin{equation}
    Y_i = \alpha + \beta_1\,\text{T1}_i + \beta_2\,\text{T2}_i + \varepsilon_i
    \label{eq:main}
  \end{equation}
  %
  where $\text{T1}_i$ indicates the allied commitment condition,
  $\text{T2}_i$ the allied restraint condition, and the 
  control is the reference category. Equation~\eqref{eq:main}
  is estimated via OLS with robust standard errors. 

  A covariate-adjusted model specification includes pre-treatment controls
  to improve precision. The same specification is estimated again using the continuous escalation variables as the dependent variable.

The same specification is estimated with the continuous measures of support for defiance as dependent variables (deescalation, sanctions, deploying ships, conducting a military strike). Following \textcite{powers_psychology_2023}, each item in the continuous
escalation battery is estimated in a separate regression using the specification in Equation~\eqref{eq:main}, with the individual item as the dependent variable. 

  %-------------------------------------------------------------------
  \subsection{Exploratory Analysis}
  %-------------------------------------------------------------------

  H3 and H4 test whether the relationship between the mechanism measures
  and support for each of the continuous defiance variables is moderated by allied commitment. Informally:
  %
  \[
    \texttt{escalation} \sim \texttt{mechanism} + \texttt{condition}
    + \texttt{mechanism} \times \texttt{condition}
  \]

  \paragraph{H3 (Reactance)}
  %
  \begin{equation}
    S_i = \alpha
        + \beta_1\,R_i
        + \beta_2\,\text{T1}_i + \beta_3\,\text{T2}_i
        + \delta_1\,(R_i \times \text{T1}_i)
        + \delta_2\,(R_i \times \text{T2}_i)
        + \varepsilon_i
    \label{eq:h3}
  \end{equation}
  %
  where $S_i$ is support for escalation and $R_i$ is the reactance
  score. 

  \paragraph{H4 (Reputation)}
  %
  \begin{equation}
    S_i = \alpha
        + \beta_1\,P_i
        + \beta_2\,\text{T1}_i + \beta_3\,\text{T2}_i
        + \delta_1\,(P_i \times \text{T1}_i)
        + \delta_2\,(P_i \times \text{T2}_i)
        + \varepsilon_i
    \label{eq:h4}
  \end{equation}
  %
  where $P_i$ is the reputational concern score. 

  %-------------------------------------------------------------------
  \subsection{Heterogeneous Effects}
  %-------------------------------------------------------------------

  Exploratory heterogeneous effects are estimated by interacting
  treatment indicators with the pre-treatment militant internationalism
  variables:
  %
  \begin{equation}
    Y_i = \alpha
        + \beta_1\,\text{T1}_i + \beta_2\,\text{T2}_i
        + \gamma\cdot\text{MilInt}_i
        + \delta_1\,(\text{T1}_i \times \text{MilInt}_i)
        + \delta_2\,(\text{T2}_i \times \text{MilInt}_i)
        + \varepsilon_i
  \end{equation}
  %
  The $\delta_k$ coefficients test whether the treatment effect of
  condition $k$ is larger among hawkish respondents. Analogous
  specifications are estimated with cooperative internationalism and
  isolationism as moderators.
%-------------------------------------------------------------------
\section{Power Analysis}
%-------------------------------------------------------------------



%-------------------------------------------------------------------
\section{Case Selection}
%-------------------------------------------------------------------

Denmark is selected as the primary case for three reasons. First, as a small NATO member with direct exposure to Russian military pressure in the Baltic Sea, Denmark exemplifies the class of states for which the theory is most relevant: states that face a materially stronger adversary and whose security depends substantively on alliance credibility. The naval scenario is ecologically valid and does not require respondents to reason about an unfamiliar threat context. Second, Denmark's multiparty political system provides sufficient within-sample variation in foreign policy orientations to test heterogeneous treatment effects. Third, studying a non-US, non-great-power sample directly addresses the sample selection bias in the existing experimental coercion literature.

The survey instrument names Denmark as the target country and Russia as the adversary. This choice prioritizes ecological validity and allows the results to speak directly to a realistic policy scenario, at the cost of some internal validity with respect to respondents' pre-existing beliefs about Russia's capabilities and intentions. Including a pre-treatment measure about respondents' attitudes towards Russia helps mitigate this concern.

\printbibliography

\end{document}
